\section{Introduction}
Since \citet{Downs57} political scientists have been asking the question: why vote? 
Voters in local school board elections may have been asking this question even 
earlier. Despite the fact that in general school boards are in charge of levying 
a substantial share of many citizens' tax bill through property taxes and 
collectively manage expenditures annually equivalent to the budget of the US 
Department of Defense, when school board elections are held off-cycle turnout 
and voter interest in them appears to be remarkably low \citep{RefWorks:333; RefWorks:321}. 
This despite the fact that the community stake in quality education is large, not 
just for parents of the students, but for the social, economic, and political 
impacts that high or low quality schools can have on a community. %citation needed

We know that voter turnout even in national on-cycle elections in the US remains 
lower than many other democracies \citep{BlaisTurnout2006}. We also have an extensive 
literature seeking to explain patterns in turnout between years and within years 
between social groups \citep{ParticAnnReview2004, econvoting2007, Galston2001, EconandOutcomes2000}.
The specific puzzle of voter turnout in school board elections has even spawned 
its own set of theories with empirical tests \citep{RefWorks:311,RefWorks:243,RefWorks:340}.

Despite this rich ground, much work remains to be done to adequately test the 
mechanisms at play in predicting turnout in off-cycle, local, non-partisan elections 
like school board elections. Furthermore, theories of participation in on-cycle and 
higher profile elections benefit from examining the limits of their generalizability and 
exploring the mechanisms that contribute to this. In this paper I seek to explore 
voter turnout in school board elections through the lenses of the major theories of 
voter turnout. I begin by reviewing these theories. Then, I move into a discussion 
of school boards and the contribution a study of school boards can bring to 
the understanding of these theories. Then I discuss the unique panel dataset of 
school board results that I have collected and the methods with which I will test 
these competing theories on the data. Finally, I conclude with an assessment of 
the performance of each of the theories and the foundations for future work. 


\section{Why Vote?}

Setting aside important though largely settled issues of access to the polls, 
enfranchisement laws, administrative and cognitive barriers to casting a ballot, and 
freedom from intimidation -- all thankfully of minimal concern in American school 
board elections -- there remains a wealth of intricate and interrelated factors 
which shape a voter's decision to cast a ballot. I broadly group this literature 
into three major strands and include the theories that have sprung up uniquely 
around school boards as a fourth category. My aim is not to fully summarize this 
rich theoretical ground, but to reframe these debates in light of the particular 
challenge of school board elections. As prologue, I succinctly summarize the 
unique case against voting in school board elections as follows:

\begin{itemize}
	\item I don't know when the election is, or I forgot. (salience)
	\item I don't know who is running. (attention)
	\item I don't have a choice - no one is challenging the incumbent. (opportunity)
	\item I don't know care which candidate wins. (lack of partisan cues)
	\item I'm satisfied with my schools. (lack of impetus)
	\item I don't care how schools perform or what the property tax is (salience)
	\item My vote won't matter
\end{itemize}

These concerns of our average potential voter are important to keep in mind as 
the major strands of voter participation theory are sketched out in the next 
sections. Macro and micro-foundations of voter participation. Why do individuals 
votes? Why do districts exhibit lower or higher turnout?

\begin{itemize}
	\item Rational Choice
	\item Socialization and Education
	\item Issue Voting
	\item Dissatisfaction and Protest
\end{itemize}

\subsection{Rational Choice}

\citet{Downs57}'s elegant statement and \citet{VoteCalculus68}'s ellaboration of the paradox 
posed by citizens voting touched off a decades long project in the discipline to understand first 
why voters would ever vote in the first place, and next why some choose to exercise 
their right to vote, and others stay home. \citet{nonvoting74} rejected the notion 
that voters were rationally assigning probabilities to different outcomes of their 
voting decision and posited instead that citizens seek to minimize their regret 
at the outcome of the election and the role their decision to vote played. However, 
these theories are significantly challenged by the fact that voting is "for many 
people, most of the time, a low-cost, low-benefit activity" \citep{aldrich1993}. 

Some of the information burden placed on voters is met by the strategic emergence 
of candidates \citep{aldrich1993}. If \citep{JacobsonBook} is correct, then the 
emergence of a strong candidate provides a clear signal to voters that the election 
is going to be close, which should drive up turnout. How this works in a low-salience 
and low-information election remains to be seen. (CAN BE TESTED BY INCLUDING EMERGENCE 
OF PRIOR CANDIDATE ON THE BALLOT AS A COVARIATE IN ANALYSIS)

Much of the long-term grounding of this school of thought is grounded in partisan 
attachment as a key supporter of turnout by reducing costs, mobilizing voters, and 
making democratic duty more clear to voters \citep{aldrich1993}. 

An important feature of this school of thought is that much of the empirical testing 
of these theories has been done in situations where voters have a high chance of being 
incidentally, that is without incurring costs, exposed to information about the election 
because of the high profile nature of these races \citep{aldrich1993}. The informational 
cues provided by partisan signals in such races even further reduces this. School board
races, for the most part, feature neither of these cost reducing features. It seems unlikely 
that these higher costs are accompanied by higher benefits in such elections. Though 
voters are assured a higher likelihood of casting the deciding vote, they are much less certain 
that the election of candidate A over candidate B will make much of a difference in 
their quality of life in the immediate or medium term. 

School boards may prove to be the electoral unit for which the Downsian hypothesis 
is the most strained. In many state school boards are sub-county administrative 
entities without partisan office holders and who often elect multiple members in 
one election cycle. Furthermore, the low level of information about school board 
candidates, their positions on issues, their likely vote share in an election, and 
the voters' own feelings on issues facing the school board make this environment 
dramatically different than the presidential election. What's more, the likelihood 
that the vote of any single voter matters in a school district is orders of 
magnitude larger than for state wide or nation wide office. 

\subsection{Socialization and Education}

\subsubsection{Socioeconomic Conditions}


Carlson
Civic knowledge and participation

\subsection{Issue Voting}

Closeness (Blas p. 119) - can we test closeness as a predictor with school boards 
where we can argue that voters have little ability to predict the closeness 
of school board elections? 


Teacher unions
Religious issues
Taxation

\subsection{Dissatisfaction and Protest}

Throw the bums out. Dahl to Alsbury. 


\section{Data and Methods}


\section{Results}


\section{Conclusion}

